\section{复利敏感性}

复利可来自资产净值变化和实际收到后再投入的现金，但本章不把两者混成固定分红假设。下式仅使用名义总回报参数$r$作税费、交易成本和基金费用之前的公式演示；没有年度追加、取现、再投资时点、通胀或市场路径。它不记录、估计或承诺$D_0$。

\[
V_{10}=\text{CNY200,000}\times(1+r)^{10}.
\]

\begin{exhibitbox}[十年总回报公式敏感性]
\small
\begin{tabularx}{\exhibitboxwidth}{P P P P}
\toprule
\textbf{年化参数} & \textbf{初始本金} & \textbf{第十年公式终值} & \textbf{公式累计增值} \\
\midrule
3.0\% & CNY200,000 & CNY268,783 & CNY68,783 \\
4.5\% & CNY200,000 & CNY310,594 & CNY110,594 \\
6.0\% & CNY200,000 & CNY358,170 & CNY158,170 \\
\bottomrule
\end{tabularx}
\end{exhibitbox}

表中数值只由公式计算，不是对十年后资产规模的承诺，也不含交易成本、税费、基金费用、年度追加、实际分配、再投资时点、通胀或市场路径。

\section{实际分配的再投资规则}

\begin{enumerate}
  \item 每季记录$D_0$的到账日、税费、费用、资产来源、净额和结算凭证；不以公告金额、历史收益率或公式参数$r$代替到账记录。
  \item 只有当资产层的资格开关为1且未超过硬上限时，实际现金才可补向该资产层。
  \item 优先用实际现金和新增资金恢复目标权重；若无合格缺口，资金留在组合内受保障流动性。
  \item 若用户将现金或本金改作消费、教育、医疗、购房或债务偿还，本报告的“全部再投资”假设即失效，必须另建用途和税后现金流模型。
\end{enumerate}

\section{分配削减敏感性}

\begin{exhibitbox}[已实现现金的审慎重算]
\small
\begin{tabularx}{\exhibitboxwidth}{L{3cm}P P}
\toprule
\textbf{情景} & \textbf{计算} & \textbf{含义} \\
\midrule
不削减 & $D_0$ & 仅反映过去12个月实际净到账现金。 \\
削减25\% & $0.75D_0$ & 可再投资现金下降，非概率判断。 \\
削减50\% & $0.50D_0$ & 不以杠杆或卖出防守资产弥补。 \\
全部归零 & $0$ & 组合风险上限不变；只是没有可再投资现金。 \\
\bottomrule
\end{tabularx}
\end{exhibitbox}

\qualitynote{本案不设“第一年分红”“十年累计分红”或固定取现额；只有产生实际现金记录后，才可衡量$D_0$及其削减敏感性。}
